% ==========================================
% CAPÍTULO: CONCLUSIONES (ENFOQUE EN MPI)
% ==========================================

\chapter{Conclusiones}

Este capítulo presenta las reflexiones finales centradas en la práctica de programación de procesos paralelos con MPI realizada en la asignatura. Se resumen los aprendizajes técnicos y conceptuales al diseñar, compilar, ejecutar y analizar programas distribuidos con MPI, así como las dificultades de implementación, las lecciones aprendidas y las recomendaciones para trabajos futuros.

\newpage
\section{Conclusión de Bustillos Cruz Jonatan}
% Escribir conclusion individual.
En el desarrollo de esta práctica, entendí la evolución del procesamiento paralelo y la importancia de MPI como una herramienta fundamental para la programación de aplicaciones distribuidas modernas. En primer lugar, hemos pasado de los primeros enfoques de paralelismo a través de hilos y procesos locales a la capacidad de ejecutar programas en múltiples nodos interconectados, lo que ha revolucionado la forma en que abordamos problemas computacionales complejos hoy en día, dado que la naturaleza suele ser más profunda y difícil de simular de lo que pensamos. En este contexto, hacen falta máquinas capaces que, al no caber en un solo chip, nos dejan sin más alternativa que comunicar miles de millones de núcleos, así como gran cantidad de RAM y almacenamiento. Por ello, herramientas como MPI se vuelven esenciales para aprovechar al máximo el potencial de estas arquitecturas distribuidas. 

Si bien al comienzo parece complejo, realmente el proceso fue diseñado para ser muy accesible. El uso de \textit{wrappers} de MPI, como \texttt{mpicc} para la compilación y \texttt{mpirun} para la ejecución, simplifica enormemente el proceso de desarrollo, permitiendo a los programadores centrarse en la lógica de su aplicación sin preocuparse por los detalles de bajo nivel de la comunicación entre procesos. En cambio, esto no sucede de la misma manera al realizar procesos de forma manual, donde es necesario gestionar explícitamente la creación de procesos, la comunicación entre ellos y la sincronización, lo que puede ser propenso a errores y requiere un conocimiento profundo del sistema operativo. En contraste, MPI abstrae estas complejidades, proporcionando una interfaz clara y consistente para el desarrollo de aplicaciones paralelas, lo que facilita la adopción de esta tecnología incluso para quienes no son expertos en sistemas distribuidos. Asimismo, esto se suma a la importancia de la estandarización y a la creación de funciones que aprovechan la configuración previa hecha con \texttt{mpirun} y permiten a los desarrolladores escribir código paralelo de manera más rápida.

A partir de aquí, solo falta aprender cada una de estas funciones; luego de familiarizarse con ellas, es mucho más sencillo escribir código realmente paralelo. Aunque por ahora no sé si será el paradigma final en el mundo de la programación paralela, pero ha demostrado ser uno de los sistemas más robustos, bien diseñados y ampliamente adoptados en la industria y la academia, lo cual hace que sea una herramienta esencial para cualquier programador que quiera trabajar en el campo del cómputo de alto rendimiento. Puede que en el futuro no use MPI directamente para crear un sistema distribuido y realizar cálculos muy complejos, pero es muy probable que se utilice como base para otras herramientas y bibliotecas. Además, me da la intuición de cómo funcionan los sistemas distribuidos, lo cual podría servirme sobre todo en el interés que tengo por el mundo de los algoritmos bioinspirados, que a menudo requieren una gran cantidad de cálculos paralelos para simular sistemas biológicos complejos y que, sin esta base, adolecen de tiempos de ejecución muy elevados.

En fin, esto no hace más que hacerme reflexionar sobre la capacidad humana de adaptarse y mejorar, lo cual es evidente en la evolución de las técnicas de paralelismo y en la creación de estándares como MPI. También, me viene a la mente que esto es lo que se denomina coloquialmente un \comillas{agujero de conejo}, es decir, un tema tan amplio y profundo que, al adentrarse en él, se descubren nuevas capas de complejidad y conocimiento. Solamente accedí a manuales básicos y libros introductorios al tema, pero, observando su potencial y amplio uso, no me queda duda de que existen muchas más funciones, técnicas y optimizaciones que se pueden aprender para mejorar el rendimiento de las aplicaciones paralelas. Además, la comunidad de MPI es muy activa, lo que significa que siempre hay nuevas herramientas, bibliotecas y recursos disponibles para quienes estén interesados en profundizar en este campo. En resumen, esta práctica me ha dado una visión general de MPI y su importancia en el mundo del cómputo paralelo, y me ha desvelado un mundo enorme que, hasta hace poco, no sabía que existía.

\newpage
\section{Conclusión de Delgado Lucero Cristian Isaac}


\newpage
\section{Conclusión de Frem Cortés José Angel}
% Escribir conclusion individual.
Durante el desarrollo de esta práctica logré comprender de manera más amplia la importancia que tiene el procesamiento paralelo dentro del área de la computación moderna, especialmente en aplicaciones donde se requiere realizar una gran cantidad de operaciones en tiempos reducidos. Antes de realizar la práctica, tenía una idea más teórica acerca del funcionamiento de MPI y del concepto de paralelismo, pero al trabajar directamente con OpenMPI y programas en lenguaje C pude entender de forma práctica cómo interactúan los procesos y cómo se distribuye el trabajo entre distintos nodos. Esto me permitió identificar la diferencia entre la programación secuencial y la programación paralela, observando que esta última permite optimizar considerablemente el uso de recursos computacionales y mejorar el rendimiento general de los programas.

Asimismo, comprendí que MPI no solamente se basa en ejecutar varios procesos al mismo tiempo, sino que también requiere una correcta coordinación y comunicación entre ellos para garantizar resultados adecuados. Durante la práctica aprendí cómo el nodo principal puede enviar información a diferentes procesos secundarios para que cada uno realice una tarea específica y posteriormente devuelva los resultados obtenidos. Este proceso me ayudó a entender la importancia de las funciones de envío y recepción de mensajes, así como la necesidad de mantener un orden correcto en la comunicación entre procesos para evitar errores durante la ejecución del programa. También observé que la sincronización es un elemento fundamental dentro del cómputo paralelo, ya que todos los procesos deben trabajar de forma coordinada para cumplir correctamente con las operaciones asignadas.

Otro aspecto importante que reforcé durante esta actividad fue el uso de herramientas y comandos relacionados con OpenMPI dentro del sistema operativo Linux. Aprendí a compilar programas paralelos utilizando \texttt{mpicc} y a ejecutar múltiples procesos mediante el comando \texttt{mpirun}, comprendiendo cómo el sistema administra los diferentes procesos generados. Además, el hecho de trabajar desde terminal permitió fortalecer mis conocimientos sobre el entorno de desarrollo en Linux y sobre la manera en que los sistemas operativos gestionan recursos relacionados con procesos y comunicación. Considero que esta experiencia fue útil porque permitió relacionar conceptos de programación, sistemas operativos y arquitectura de computadoras dentro de una misma práctica.

De igual manera, esta práctica me permitió desarrollar habilidades relacionadas con el razonamiento lógico y la resolución de problemas. Al trabajar con procesos distribuidos fue necesario analizar cuidadosamente cómo dividir las tareas, cómo enviar correctamente la información y cómo recuperar los resultados finales desde el nodo principal. Esto me hizo comprender que la programación paralela requiere mayor planeación y organización en comparación con la programación secuencial, debido a que cualquier error pequeño puede afectar la ejecución completa del sistema. También aprendí que el diseño de programas paralelos debe considerar aspectos como el balance de carga y la reducción de comunicación innecesaria entre procesos para obtener un mejor rendimiento.

Además, logré identificar algunas de las ventajas y limitaciones que presenta MPI dentro del desarrollo de aplicaciones paralelas. Por una parte, observé que permite aprovechar mejor el hardware disponible y reducir tiempos de procesamiento al ejecutar varias tareas simultáneamente. Sin embargo, también comprendí que la complejidad del código aumenta debido a la comunicación y sincronización entre procesos, lo cual puede dificultar la depuración de errores y el mantenimiento del programa. A pesar de ello, considero que MPI continúa siendo una herramienta muy importante dentro del área del cómputo de alto rendimiento debido a su capacidad para trabajar en sistemas distribuidos y su amplia utilización en aplicaciones científicas, académicas e industriales.
